\documentclass[12pt]{article}
\usepackage{amsmath,amssymb}
\newcommand{\Dil}{\operatorname{Dil}}
\newcommand{\Vol}{\operatorname{Vol}}
\newcommand{\Fill}{\operatorname{Fill}}
\newcommand{\Mass}{\operatorname{Mass}}
\begin{document}
Scratchpad for the rectangle 2-dilation problem.

Interpretation: rectangles are oriented Euclidean boxes; a map
\(f:(R,\partial R)\to(S,\partial S)\) of degree 1 means the induced map on
\(H_4(\cdot,\partial\cdot;\mathbb Z)\) sends \([R]\) to \([S]\).
The constant called \(k\) in the problem is a small universal constant;
in drafts it is denoted \(\kappa\) to avoid conflict with \(k\)-dilation.

Important correction after referee report.  The earlier proposed theorem
\[
 M_2/(S_3S_4)+M_3/(S_1S_3S_4)
 +M_4/(S_1S_2^{1/2}S_3^{3/2}S_4)\ge c
\]
for arbitrary two-parameter sweepouts is false.  Counterexample:
\(S=[0,1]^2\times[0,L]^2\) and cycles
\([0,1]^2\times\{u\}\times\{v\}\).  Then \(M_2=1,M_3=L,M_4=L^2\), so the
left side is \(L^{-2}+L^{-1}+L^{-1/2}\to0\).  Any valid proof must use
more of the map \(f\) than the fiber sweepout: specifically, the
\(\Dil_2\) bound controls mixed two-faces of a cubical decomposition of
\(R\).

Precise Guth estimate from arXiv:0710.0403, Estimate 1.  Let
\(Q_i=S_i/R_i\).  If \(\Phi:(U,\partial U)\to(S,\partial S)\), \(U\subset R\),
has nonzero degree, then for \(0\le j<k\le l\le n\),
\[
 \operatorname{dil}_k(\Phi)\ge c(n)Q_1\cdots Q_j
   (Q_{j+1}\cdots Q_l)^{(k-j)/(l-j)}.
\]
Equivalently, if \(\Phi\) is \(k\)-contracting and \(0<j<k<l\),
\[
 \left({R_1\cdots R_j\over S_1\cdots S_j}\right)^{(l-k)/(k-j)}
 R_1\cdots R_l \ge c(n)S_1\cdots S_l .
\]
Estimate 2 includes a factor \(D^{(k-j)/(n-j)}\) but for degree \(D=1\)
coincides with the \(l=n\) case.

Consequences for \(n=4,k=2\), plus \(\Dil_3,\Dil_4\le1\) from singular
values:
\[
R_1R_2\gtrsim S_1S_2,\quad R_1R_2R_3\gtrsim S_1S_2S_3,\quad
V_R\gtrsim V_S,
\]
\[
R_1^2R_2R_3\gtrsim S_1^2S_2S_3,\quad
R_1^2V_R\gtrsim S_1^2V_S,
\]
and from \(k=3,j=2,l=4\),
\[
(R_1R_2)^2R_3R_4\gtrsim (S_1S_2)^2S_3S_4.
\]
The \(k=2,j=1,l=4\) inequality gives, if \(R_1\le\kappa S_1\),
\[
V_R\gtrsim \kappa^{-2}S_1S_2S_3S_4.
\]
This implies the desired volume lower bound
\(V_R\gtrsim \kappa S_1S_2^{1/2}S_3^{3/2}S_4\) only when
\(S_3/S_2\lesssim \kappa^{-6}\).  The difficult range is
\(S_3/S_2\gg \kappa^{-6}\).

Scaling test showing Guth monomial estimates alone are insufficient:
take
\[
S=(1,1,\varepsilon^{-6},\varepsilon^{-6}),\quad
R=(\varepsilon,\varepsilon^{-4},\varepsilon^{-5},\varepsilon^{-6}).
\]
Then \(R_1=\varepsilon S_1\), \(R_3R_4=\varepsilon S_3S_4\), and
\[
V_R=\varepsilon\, S_1S_2^{1/2}S_3^{3/2}S_4.
\]
The quotients are \(Q=(\varepsilon^{-1},\varepsilon^4,\varepsilon^{-1},1)\).
All Estimate 1 monomials for \(k=2\),
\[
Q_1Q_2,\ (Q_1Q_2Q_3)^{2/3},\ (Q_1Q_2Q_3Q_4)^{1/2},\
Q_1(Q_2Q_3)^{1/2},\ Q_1(Q_2Q_3Q_4)^{1/3},
\]
are \(O(1)\) or smaller.  The \(k=3\) monomial
\(Q_1Q_2(Q_3Q_4)^{1/2}\) is also small.  Thus the desired result needs a
genuinely mixed-face strengthening, not just the known monomial theorem.

A plausible sufficient mixed-face estimate is
\[
 {R_3R_4\over S_3S_4}
 +{R_1\sqrt{R_3R_4}\over S_1\sqrt{S_3S_4}}
 +{V_R\over S_1S_2^{1/2}S_3^{3/2}S_4}\ge c.
\]
Under \(R_1\le\kappa S_1\) and \(R_3R_4\le\kappa S_3S_4\), the first two
terms are \(\le\kappa\) and \(\le\kappa^{3/2}\), yielding the desired
volume lower bound.  Need prove or locate this estimate.  It should come
from an anisotropic version of Guth's tightening construction where the
initial \(2\)-skeleton records all mixed \(2\)-faces.  The geometric mean
\(R_1\sqrt{R_3R_4}\) is controlled even though the individual mixed face
areas \(R_1R_3\), \(R_1R_4\) need not both be small when \(S_4/S_3\) is
large.

Ideas tried:
1. Packing-width estimates for \(2\)-width and \(3\)-width of preimages of
target subrectangles reproduce standard Guth monomial lower bounds but so
far do not generate the factor \((S_3/S_2)^{1/2}\).
2. Applying Estimate 1 to \(\Dil_3\le1\) gives
\((R_1R_2)^2R_3R_4\gtrsim(S_1S_2)^2S_3S_4\), but combining with
\(R_3R_4\le\kappa S_3S_4\) gives only \(R_1R_2\gtrsim\kappa^{-1/2}S_1S_2\).
3. Coarea on \((f_3,f_4)\) gives the basic volume lower bound
\(V_R\gtrsim S_1S_2S_3S_4\), but does not use \(R_1\) small.
4. Product constructions from 3D snake maps fail after crossing with an
interval because the mixed \(2\)-dilation involves the large \(1\)-dilation
of the snake map.  This supports the idea that mixed faces are the missing
obstruction.



New direction after receiving council/compute feedback: the proposed theorem may be
false for degree-one maps (as opposed to diffeomorphisms or maps with additional
rigidity).  Guth's Chapter 8.2 pinching maps appear to realize the scaling test as an
actual bounded-2-dilation map.

Asymptotic counterexample in length parameter \(L=\varepsilon^{-1}\):
\[
 R^0=(L^{-1},L^4,L^5,L^6),\qquad S=(1,1,L^6,L^6).
\]
The desired alternatives fail at scale \(1/L\):
\[
 R_1^0/S_1=L^{-1},\quad
 R_3^0R_4^0/(S_3S_4)=L^{11}/L^{12}=L^{-1},
\]
\[
 \operatorname{Vol}(R^0)/(S_1S_2^{1/2}S_3^{3/2}S_4)=L^{14}/L^{15}=L^{-1}.
\]

Guth map sequence (all degree one maps of pairs with \(\Dil_2\le C\), constants
independent of \(L\)):
1. ``Stretch the short side''/technical remark with stretch \(A=L\):
\[
(L^{-1},L^4,L^5,L^6)\to (1,L,L^5,L^6).
\]
This uses \(A^4=L^4\le R_2/R_1=L^5\).
2. Pinching map with \(a=1\), \(b=L^{12}\):
\[
(1,L,L^5,L^6)\to(1,1,1,L^{12}),
\]
using \(a\le R_1\) and \(a^2b=L^{12}=R_2R_3R_4\).
3. Last-factor splitting/technical remark:
\[
(1,1,1,L^{12})\to(1,1,L^6,L^6).
\]
This is the non-linear step that folds the long \(L^{12}\) direction into two
\(L^6\) directions while controlling mixed 2-dilation by pinching the exceptional
set.  It is important not to replace it by a product of the identity with a planar
snake map, because that would have large mixed \(2\)-dilation.

The composition has \(\Dil_2\le C_0\).  Scaling the domain by
\(\lambda=C_0^{1/2}\) and precomposing by the homothety gives a map
\(\lambda R^0\to S\) with \(\Dil_2\le1\).  The three bad ratios become
\(\lambda/L,\lambda^2/L,\lambda^4/L\), so for any proposed \(\kappa>0\) choosing
\(L\gg_\lambda \kappa^{-1}\) violates the conclusion.

Need verify exact statement of Guth technical remark if challenged: Chapter 8.2,
map 3 and its technical remark give the short-side stretch and splitting variants;
map 4 gives pinching.  The current answer cites these as standard constructions and
recalls the retraction/pinching idea.  A future refinement could reproduce the
construction in full detail or quote exact page numbers.


Correction/refinement of the counterexample sequence: the earlier tentative
sequence using
\((L^{-1},L^4,L^5,L^6)\to(1,L,L^{11/2},L^{11/2})\to(1,1,1,L^{12})\to(1,1,L^6,L^6)\)
used an overgeneralized version of Guth's technical remark.  The exact technical
remark ties the stretch parameter \(A\) to the new third side in the interpolated
target \(AR_1\times A^{-3}R_2\times A\times B\).

A clean valid exponent choice using the exact statements is:
\[
R^0=(L^{-7},1,L,L),\qquad
S=(L^{-13/2},L^{-13/2},1,L^{11/2}).
\]
Map sequence:
\[
(L^{-7},1,L,L)
 \xrightarrow[\text{short-side stretch }A=L]{}
(L^{-6},L^{-3},L,L),
\]
because \(L^4\le L^7=R_2/R_1\).
Then pinching with \(a=L^{-13/2}\), \(b=L^{12}\):
\[
(L^{-6},L^{-3},L,L)
 \to (L^{-13/2},L^{-13/2},L^{-13/2},L^{12}),
\]
because \(a\le L^{-6}\) and \(a^2b=L^{-1}=L^{-3}LL\).  Guth states some inequalities strictly; equalities here can be perturbed by fixed factors and do not affect asymptotics.
Finally the interpolated short-side/stretch-split map with \(A=1\):
\[
(L^{-13/2},L^{-13/2},L^{-13/2},L^{12})
 \to (L^{-13/2},L^{-13/2},1,L^{11/2}),
\]
because \(A^4=1\le R_2/R_1=1\) and
\(R_3=L^{-13/2}\le1\le(R_3R_4)^{1/2}=L^{11/4}\).

The bad ratios for \(R^0,S\):
\[
R_1/S_1=L^{-1/2},\quad
R_3R_4/(S_3S_4)=L^2/L^{11/2}=L^{-7/2},
\]
\[
\operatorname{Vol}(R^0)/(S_1S_2^{1/2}S_3^{3/2}S_4)
=L^{-5}/(L^{-13/2}L^{-13/4}L^{11/2})=L^{-3/4}.
\]
After scaling the domain by \(\lambda=C_0^{1/2}\), these become
\(\lambda L^{-1/2}\), \(\lambda^2L^{-7/2}\), \(\lambda^4L^{-3/4}\).  Hence they
all tend to zero.  This gives a genuine counterexample if the cited Guth maps are
accepted.



Round 3 update / current mathematical state.

The round-2 counterexample using the map
\[
(L^{-13/2},L^{-13/2},L^{-13/2},L^{12})\to
(L^{-13/2},L^{-13/2},1,L^{11/2})
\]
is invalid.  Guth's published Estimate 1 for degree-one maps of pairs applies with
\(j=0,k=2,l=3\) and gives a lower bound \(\Dil_2\gtrsim L^{13/3}\).  Therefore the
``interpolated technical remark'' from the thesis cannot be used in that broad form;
at minimum it has hidden hypotheses or was superseded/corrected by the later
published estimate.

A rigorous partial result that remains useful:
For degree-one \(f:R\to S\) with \(\Dil_2(f)\le1\), if
\[
 C R_1R_2^2 \le S_1S_2S_3,
\]
then
\[
 R_3R_4\ge cS_3S_4.
\]
Proof: pull back the central relative 2-planes
\([0,S_1]\times[0,S_2]\times\{y_3,y_4\}\).  Their filling volume in \(S\) is
\(\gtrsim S_1S_2S_3\).  Since \(\Dil_3(f)\le1\), their preimages have at least this
filling volume in \(R\).  Guth's small relative 2-cycle filling lemma says any
relative 2-cycle of area \(<cR_1R_2\) fills with volume \(\lesssim R_1R_2^2\).  Hence
the pulled-back cycles have area \(\gtrsim R_1R_2\).  Coarea for
\((x_3,x_4)\circ f\) yields \(\Vol(R)\gtrsim R_1R_2S_3S_4\), hence the conclusion.

Thus the only remaining case (assuming the first desired alternative is false) is
\[
 R_1R_2^2\gtrsim S_1S_2S_3,\qquad R_1\le\kappa S_1.
\]
For \(\kappa\) small this implies \(R_2^2\gg S_2S_3\).

Naive attempt for the remaining case:
Pull back the intervals
\[
\{y_1\}\times[0,S_2]\times\{y_3\}\times\{y_4\}
\]
and try to use local isoperimetry to show their preimages have length
\(\gtrsim (S_2S_3)^{1/2}\); coarea for \((x_1,x_3,x_4)\circ f\) would then give
\[
 \Vol(R)\gtrsim S_1S_3S_4(S_2S_3)^{1/2}
 =S_1S_2^{1/2}S_3^{3/2}S_4.
\]
Obstacle: those intervals are relative 1-cycles in \(S\) that can fill through an
\(x_1\)-face with area \(O(S_1S_2)\), not \(\gtrsim S_2S_3\).  Need a genuinely mixed
family of cycles or a chain-level tightening that prevents escape through the short
\(x_1\) direction.

Relevant source details from Guth thesis:
- Proposition 8.1.2: a relative 2-cycle in a rectangle of area \(<\frac12 R_1R_2\)
  fills with volume \(\lesssim R_1R_2^2\) (OCR sometimes drops the square).
- Proposition 8.1.3: for a 2-contracting diffeomorphism, if
  \(CR_1R_2^2<S_1S_2S_3\), then \(R_3R_4\gtrsim S_3S_4\).  The proof adapts to
  degree-one maps of pairs using slicing currents.
- Page 130 rational homotopy inequalities for diffeomorphisms in dimension 4 include
  boundary \(R_2R_3^2R_4\gtrsim S_2S_3^2S_4\) and double inequalities
  \(R_1R_2R_3^2R_4\gtrsim S_1S_2S_3^2S_4\),
  \(R_1R_2^2R_3^2R_4\gtrsim S_1S_2^2S_3^2S_4\).  These do not by themselves imply
  the problem's desired estimate; the scaling test \(S=(1,1,L^6,L^6)\),
  \(R=(L^{-1},L^4,L^5,L^6)\) satisfies them but still has the bad ratios \(L^{-1}\).

Potential counterexample route still unresolved:
The scaling test would disprove the problem if one could build a bounded
2-dilation degree-one map.  A tempting decomposition is
\[
(L^{-1},L^4,L^5,L^6)\to(1,L,L^5,L^6)\to(1,1,L^6,L^6).
\]
The first arrow is Guth's short-side-stretch technical map with parameter \(A=L\)
and is not contradicted by Estimate 1.  The second ``middle-side snake''
\((1,L,L^5,L^6)\to(1,1,L^6,L^6)\) is not one of the listed safe Guth maps.  It is not
ruled out by the interior published Estimate 1, but it is ruled out for diffeomorphisms
by the boundary rational-homotopy inequality
\(R_2R_3^2R_4\gtrsim S_2S_3^2S_4\) (left \(L^{17}\), right \(L^{18}\)).  For degree-one
maps of pairs it remains unclear whether such a map can exist; if the boundary
inequality applies to arbitrary boundary degree-one maps with bounded 2-dilation,
then the residual middle-side snake is impossible.



Technical lemmas for the slicing approach.

I made the first-alternative argument fully current/chain based.  For a central
target plane
\[
P_y=[0,S_1]\times[0,S_2]\times\{y_3\}\times\{y_4\}
\]
a relative filling \(Y\) in \(S\) has volume \(\gtrsim S_1S_2S_3\).  The clean
proof uses the calibration
\[
\omega=dx_1\wedge dx_2\wedge d\psi,
\]
where \(\psi\) is a 1-Lipschitz function on the \((x_3,x_4)\)-rectangle, zero on
its boundary and \(\psi(y)\gtrsim S_3\) for central \(y\).  Boundary terms vanish:
on \(x_1,x_2\)-faces \(dx_1\wedge dx_2=0\), and on \(x_3,x_4\)-faces
\(\psi=0\).  This also fixes the previous objection that a filling might escape
through \(x_4\).

Slicing naturality: for \(F=(x_3,x_4)\circ f\) and a.e. regular \(y\),
\[
f_\#\langle [R,\partial R],F,y\rangle
=\langle f_\#[R,\partial R],(x_3,x_4),y\rangle
=P_y
\]
as relative currents.  Hence any filling of the slice in \(R\) pushes forward,
using \(\Dil_3(f)\le1\), to a filling of \(P_y\).

Guth's rectangular isoperimetric profile for \(k=2,n=4\): if a relative
2-cycle has area \(A\le cR_1R_2\), then
\[
I_R^2(A)\le C\max(A^{3/2},A^2/R_1),
\]
and in particular \(I_R^2(A)\le CR_1R_2^2\).  This gives the rigorous first
alternative:
\[
CR_1R_2^2\le S_1S_2S_3\quad\Longrightarrow\quad R_3R_4\gtrsim S_3S_4.
\]

A useful but insufficient profile consequence for the complementary case:
For the central slices \(Z_y\),
\[
\Mass Z_y\gtrsim
\min\{R_1R_2,\ (S_1S_2S_3)^{2/3},\ (R_1S_1S_2S_3)^{1/2}\}.
\]
Coarea over \(y_3,y_4\) gives the same expression times \(S_3S_4\) as a lower
bound for \(\Vol(R)\).  The term \((S_1S_2S_3)^{2/3}\) would imply the desired
fiber-area lower bound \(S_1(S_2S_3)^{1/2}\), since
\((S_1S_2S_3)^{2/3}/(S_1S_2^{1/2}S_3^{1/2})
=(S_2S_3/S_1^2)^{1/6}\ge1\).  The obstruction is the alternative
\((R_1S_1S_2S_3)^{1/2}\), which is smaller by \((R_1/S_1)^{1/2}\).

Guth Estimate 1 with \(j=1,l=4,k=2\) gives
\[
R_1^3R_2R_3R_4\gtrsim S_1^3S_2S_3S_4.
\]
Under \(R_1\le\kappa S_1\),
\[
\Vol(R)\gtrsim \kappa^{-2}S_1S_2S_3S_4,
\]
which proves the desired volume bound when \(S_3/S_2\lesssim \kappa^{-6}\).
Thus the hard range is \(S_3/S_2\gg_\kappa1\), and, after the first alternative
reduction, \(R_1R_2^2\gtrsim S_1S_2S_3\).

Potential next attack: upgrade the profile consequence by using the degree of
the map \(Z_y\to P_y\), not merely the area and filling volume of \(Z_y\).  A
formal slicing of \(Z_y\) by the target \(x_1\)-coordinate suggests the desired
estimate \( \Fill(Z_y)\lesssim A^2/S_1\), which would imply
\(A\gtrsim S_1(S_2S_3)^{1/2}\).  The obstacle is that \(f_1\) is not
1-Lipschitz; however on \(Z_y\) the Jacobian bound
\(|df_1\wedge df_2|\le1\) and the degree-one condition to
\([0,S_1]\times[0,S_2]\) may provide the missing weighted coarea control.
This is the most concrete remaining subproblem.




\section*{Single-fiber no-go and family-level obstruction}

A proposed route was to improve the source filling upper bound for each central
fiber \(Z_y=f^{-1}([0,S_1]\times[0,S_2]\times\{y_3,y_4\})\) from the Guth-profile
branch \(A^2/R_1\) to \(A^2/S_1\).  This is false under slice-only hypotheses.
Model: choose a relative rectangle
\[
Z=[0,r]\times[0,L]\times\{(h,h)\}\subset R,\qquad A=rL,
\]
and map it to \([0,S_1]\times[0,S_2]\) by
\[
(u,v)\mapsto(S_1u/r,S_2v/L).
\]
The Jacobian is \(S_1S_2/A\le1\) if \(A\ge S_1S_2\), and the map has relative
degree one.  A calibration \(\omega=dx_1\wedge dx_2\wedge d\psi\), with
\(\psi(h,h)\sim h\), shows that its filling volume in the ambient rectangle is
\(\gtrsim hA\).  Taking \(h=S_1S_2S_3/A\) gives filling volume
\(\gtrsim S_1S_2S_3\) at area
\[
A\sim (rS_1S_2S_3)^{1/2},
\]
which is exactly the offending \(A^2/R_1\) branch when \(r=R_1\).  Hence a proof
of the problem must use the fact that the \(Z_y\)'s are fibers of a single
four-dimensional map, not merely the degree/Jacobian data on each fiber.

Current rigorous partial results:
\[
CR_1R_2^2\le S_1S_2S_3 \Longrightarrow R_3R_4\gtrsim S_3S_4.
\]
If the first desired alternative fails (for small \(\kappa\)), then
\[
R_1R_2^2\gtrsim S_1S_2S_3,\qquad
R_2^2\gtrsim \kappa^{-1}S_2S_3.
\]
Guth's Estimate 1 with \(j=1,k=2,l=4\) gives
\[
R_1^2\operatorname{Vol}(R)\gtrsim S_1^3S_2S_3S_4,
\]
so \(R_1\le\kappa S_1\) implies
\[
\operatorname{Vol}(R)\gtrsim \kappa^{-2}S_1S_2S_3S_4.
\]
This proves the desired second alternative in the range
\(S_3/S_2\lesssim\kappa^{-6}\).  The remaining hard range is
\(S_3/S_2\gg_\kappa1\), with \(R_2^2\gg S_2S_3\), and requires a genuinely mixed
family-level estimate.

The scaling stress test remains
\[
S=(1,1,L^6,L^6),\qquad R=(L^{-1},L^4,L^5,L^6).
\]
It saturates Guth's known monomial estimates and the boundary \(B3\)-type
ellipsoid estimate, while violating the desired alternatives by a factor \(L^{-1}\).
No valid bounded-\(2\)-dilation degree-one map \(R\to S\) is known.  The tempting
route
\[
(L^{-1},L^4,L^5,L^6)\to(1,L,L^5,L^6)\to(1,1,L^6,L^6)
\]
fails at the second arrow: the boundary/ellipsoid obstruction gives
\(R_3R_4\gtrsim S_3S_4\) in the relevant case, but \(L^{11}<L^{12}\).



\section*{Round 6 boundary-trace reduction}

New potentially useful lemma.  Let \(Z\) be a relative \(2\)-cycle in a
rectangle \(R\), \(A=\Mass Z\).  For a coordinate \(x_i\), let \(L_i\) be the
mass of the trace of \(\partial Z\) on the two faces \(x_i=0,R_i\).  Then
\[
\Fill_R(Z)\lesssim R_iA+A L_i.
\]
Proof idea:
1. General relative FF after doubling gives \(\Fill_R(Z)\lesssim A^{3/2}\).
By AM-GM, \(A^{3/2}\le (R_iA+A^2/R_i)/2\).
2. Fill the trace \(\Gamma_i\) in the two \(x_i\)-faces by relative
2-chains of area \(\lesssim L_i^2\).  Adding these face chains changes \(Z\)
only by a relative-zero chain and gives a representative \(Z'\) whose boundary
does not lie on the \(x_i\)-faces.  Translate \(Z'\) linearly to \(x_i=0\);
the remaining boundary track stays in \(\partial R\).  This fills with volume
\(\lesssim R_i(A+L_i^2)\).
Taking the minimum of the \(A^2/R_i\) and \(R_iL_i^2\) terms gives
\(\lesssim R_iA+AL_i\).

For central fibers \(Z_y=\langle [R,\partial R],(f_3,f_4),y\rangle\), with
\(A_y=\Mass Z_y\) and \(L_{i,y}\) the trace on \(x_i\)-faces, the target
calibration gives \(\Fill_R(Z_y)\gtrsim F:=S_1S_2S_3\).  Hence
\[
A_y(R_i+L_{i,y})\gtrsim F.
\]
Coarea on the two \(x_i\)-faces gives
\[
\int_Q L_{i,y}\,dy\lesssim \Vol(R)/R_i,
\]
because \(J_2((f_3,f_4)|_{\{x_i=0,R_i\}})\le1\).  Jensen/optimization:
for \(q=|Q|\sim S_3S_4\),
\[
\int_Q A_y\,dy\gtrsim {Fq\over R_i+\Vol(R)/(R_iq)}.
\]
Since \(\int_Q A_y\le\Vol(R)\),
\[
R_i\Vol(R)+{\Vol(R)^2\over R_iS_3S_4}\gtrsim S_1S_2S_3^2S_4
\quad\text{for all }i.
\]
This is a genuine family-level strengthening.  It recovers, for the stress-test
dimensions \(S=(1,1,L^6,L^6)\), \(R=(L^{-1},L^4,L^5,L^6)\), an obstruction
from \(i=3\) if the constants were sharp enough, but the inequality as used
still needs to be combined with \(R_3R_4\le\kappa S_3S_4\), side ordering, and
Guth's monomial bounds.  Algebraic optimization of only
\(\Vol\gtrsim S_3S_4\sqrt{R_1S_1S_2S_3}\) and
\(\Vol\gtrsim S_1S_2S_3S_4/(R_1/S_1)^2\) leaves a gap when
\(S_3/S_2\) is extremely large, so the face inequalities for \(i=2,3,4\) or an
additional boundary/slicing input are likely needed.

Questions to resolve next:
- Can the four trace inequalities plus \(R_3R_4\le\kappa S_3S_4\) algebraically
imply \(\Vol\gtrsim \kappa S_1S_2^{1/2}S_3^{3/2}S_4\)?
- If not, is there a sharper coarea control for the trace term, e.g. involving
\(\int A_yL_{i,y}\) or simultaneous traces on both \(x_3,x_4\) faces?
- Check whether Guth's thesis contains an exact lower bound or an explicit
snake-map counterexample with the scaling
\((L^{-1},L^4,L^5,L^6)\to(1,1,L^6,L^6)\).




\section*{Round 7 correction: trace lemma is false; weighted coarea replacement}

The Round 6 one-face trace filling lemma
\[
 \Fill_R(Z)\lesssim R_i\Mass(Z)+\Mass(Z)L_i
\]
is false for arbitrary relative \(2\)-cycles.  Counterexample:
\(R=[0,1]^2\times[0,L]^2\), \(Z=[0,1]^2\times\{(L/2,L/2)\}\).  Then
\(\Mass Z=1\), the trace on (say) the \(x_1\)-faces has length \(O(1)\), but the
calibration \(dx_1\wedge dx_2\wedge d\psi(x_3,x_4)\), with
\(\psi=0\) on the \((x_3,x_4)\)-boundary and \(\psi(L/2,L/2)\sim L\), gives
\(\Fill_R(Z)\gtrsim L\).  The false proof failed in two places: Euclidean
\(\Fill\lesssim A^{3/2}\) is not true for relative cycles in long rectangles, and
a short trace in a long face need not fill inside the face with area \(O(L_i^2)\)
independent of the distance to the boundary.  Do not use inequalities (8)--(11)
from the old draft.

A valid family-level replacement is the following weighted coarea observation.
Let \(F=(f_3,f_4)\), \(Z_y=F^{-1}(y)\), and
\[
 \lambda(x)=\|d(f_1,f_2)|_{T_xZ_y}\|.
\]
At a regular point, if \(v\in T_xZ_y\) realizes \(\lambda\) and \(n\) is a unit normal
to \(Z_y\), then \(dF(v)=0\) and
\[
 |df(v)\wedge df(n)|\ge \lambda |dF(n)|.
\]
Since \(\Dil_2(f)\le1\), \(|dF(n)|\le\lambda^{-1}\), so the normal \(2\)-Jacobian
\(J_F\le C\lambda^{-2}\).  Coarea gives
\[
 \Vol(R)\ge c\int_Q\int_{Z_y}\lambda^2\,d\mathcal H^2\,dy.
\]
For each central \(y\), \((f_1,f_2):Z_y\to[0,S_1]\times[0,S_2]\) has relative
degree \(1\), so \(\int_{Z_y}|df_1\wedge df_2|\ge S_1S_2\) and
\(\int_{Z_y}\lambda^2\gtrsim S_1S_2\).  This recovers only
\(\Vol(R)\gtrsim S_1S_2S_3S_4\).  A sufficient missing lemma would be a
fiber-family energy lower bound
\[
 \int_{Z_y}\lambda^2\gtrsim S_1S_2^{1/2}S_3^{1/2}
\]
for a positive-measure set of central \(y\), using the filling lower bound
\(\Fill_R(Z_y)\gtrsim S_1S_2S_3\) and the fact that the \(Z_y\)'s are fibers of the
same ambient map.

Algebraic insufficiency of known single-fiber/profile estimates persists.  With
\(S=(1,1,L,L)\), the exponent model
\[
 r=(-1/5,\,3/5,\,1,\,1)
\]
has \(R_1R_2^2\) exactly at the first-alternative threshold and volume exponent
\(12/5\), below the desired \(5/2\).  If the last two sides are multiplied by a
small fixed constant this also models \(R_3R_4\le\kappa S_3S_4\) at the exponent
level.  It is the short-side-stretch scaling; Guth's short-side-stretch map realizes
the comparable case \(R_3R_4\sim S_3S_4\), so the missing argument must use the
strict smallness of the long-area ratio, not just monomial degree estimates.

Correct OCR for Guth thesis double pinching map (p.148): the target is
\[
 R_1^2/A \times A\times A\times B,
\]
not \(R_1/A\times A\times A\times B\).  Conditions:
\(R_1<A\), \(A^2<R_2R_3\), \(A^3B<R_1R_2R_3R_4\).  Earlier exponent searches
using \(R_1/A\) produce spurious counterexamples and must be ignored.  Random
searches with the corrected \(R_1^2/A\) output found no Guth-map counterexample
to the stated theorem; the best scalings saturate the first alternative rather than
violating it.




Further technical consolidation.

Correct weighted coarea identity.  Let \(F=(f_3,f_4)\), \(G=(f_1,f_2)\), and at
rank-two points of \(F\) let
\[
\lambda(x)=\|dG|_{\ker dF_x}\|.
\]
Then the pointwise inequality is
\[
 \lambda(x)^2 J_2F(x)\le 1.
\]
Proof: if \(T=\ker dF_x\), \(N=T^\perp\), \(v\in T\) is unit with
\(|dG(v)|=\lambda\), and \(n\in N\) is unit, then
\[
df(v)=(dG(v),0),\quad df(n)=(dG(n),dF(n)),
\]
so \(|df(v)\wedge df(n)|\ge \lambda |dF(n)|\).  Since \(\Dil_2 f\le1\),
\(\|dF|_N\|\le\lambda^{-1}\), hence \(J_2F\le\lambda^{-2}\).  Coarea gives
\[
\int_Q\int_{\langle R,F,y\rangle}\lambda^2\,d{\cal H}^2\,dy
=\int_{F^{-1}(Q)}\lambda^2J_2F\,dx\le \Vol(R).
\]
This avoids dividing by \(J_2F\); critical points simply contribute zero to the
coarea integrand.  For a.e. central \(y\), slicing naturality gives
\(f_\#\langle [R,\partial R],F,y\rangle=P_y\), so \(G\) has relative degree one on
the slice and
\[
\int_{Z_y}\lambda^2\ge \int_{Z_y}|df_1\wedge df_2|\ge S_1S_2.
\]
This recovers only \(V_R\gtrsim S_1S_2S_3S_4\).

Averaged filling estimate that would close the theorem.  Put
\[
q=|Q|\sim S_3S_4,\quad F_0=S_1S_2S_3,\quad
I=\int_Q\int_{Z_y}\lambda^2 .
\]
A sufficient family-level estimate is
\[
\int_Q \Fill_R(Z_y)\,dy
\le C\left(R_1I+\frac{I^2}{S_1q}
+\frac{R_3R_4}{S_3S_4}F_0q\right). \tag{AF}
\]
Indeed \(\Fill_R(Z_y)\gtrsim F_0\) for a.e. central \(y\), so if
\(R_3R_4\le\kappa S_3S_4\), the last term is absorbed for small \(\kappa\).  If also
\(I\le\eta qS_1(S_2S_3)^{1/2}\), then
\[
R_1I\le \eta\kappa F_0q,\qquad I^2/(S_1q)\le\eta^2F_0q
\]
(using \(R_1\le\kappa S_1\) and \(S_1\le S_2\le S_3\)), a contradiction for small
\(\eta,\kappa\).  Hence \(I\gtrsim qS_1(S_2S_3)^{1/2}\), and weighted coarea gives
\[
\Vol(R)\gtrsim S_3S_4\,S_1(S_2S_3)^{1/2}
=S_1S_2^{1/2}S_3^{3/2}S_4.
\]
The defect term in (AF) is necessary: if \(R=S\) and \(Z_y=[0,S_1]\times[0,S_2]\times\{y\}\), the family is a product foliation with
large filling but no extra weighted energy; its obstruction is exactly
\(R_3R_4/S_3S_4\sim1\).

Why simple variants fail.  A fiberwise inequality
\[
\int_{Z_y}\lambda^2\gtrsim S_1(S_2S_3)^{1/2}
\]
is false without ambient information.  In \(S=(1,1,L,L)\), a sheet spanning long
source directions and mapping them with degree one to the \((f_1,f_2)\)-square can
have \(\int\lambda^2\sim1\) while its filling in the source is very large.  Such a
single-fiber model violates the full ambient \(2\)-dilation bounds when one tries to
make it part of a family covering the \((f_3,f_4)\)-rectangle; the missing argument
must exploit the pointwise mixed bound \(\lambda\|dF|_N\|\le1\) together with the
limited source parameter area \(R_3R_4\).

Boundary and level-set ideas.  Guth's ellipsoid boundary estimate B3a says, in the
case \(R_2^2\gtrsim S_2S_3\), that a \(2\)-contracting boundary diffeomorphism
forces
\[
R_2R_3R_4\gtrsim S_2^{1/2}S_3^{3/2}S_4.
\]
The complementary first-alternative range and \(R_1\le\kappa S_1\) imply
\(R_2^2\gtrsim_\kappa S_2S_3\), so this boundary estimate has the correct
\(S_2,S_3,S_4\) pattern.  However multiplying it by \(R_1\) gives only
\(V_R\gtrsim R_1S_2^{1/2}S_3^{3/2}S_4\), not the desired \(S_1\)-factor.  A level-set
approach using \(f_1=t\) also runs into the fact that target intervals/circles at
interior \(x_1=t\) can fill through the short \(x_1\)-direction with area \(O(S_1S_2)\),
so it does not reproduce the boundary filling area \(S_2S_3\).  These failures point
back to an averaged double/parametric version of the boundary argument, rather than
a direct boundary or single-level argument.

Short-side-stretch counterexample checks.  The tempting map
\[
(L^{-1/5},L^{3/5},\epsilon L,\epsilon L)\to
(1,1,\epsilon L^{6/5},\epsilon L^{6/5})
\]
with stretch \(A=L^{1/5}\) is not one of Guth's valid short-side-stretch maps and is
ruled out by published Estimate 1: for the proposed target the quotients are
\((L^{1/5},L^{-3/5},L^{1/5},L^{1/5})\), and the \(j=1,k=2,l=4\) lower bound gives
\(\Dil_2\gtrsim L^{2/15}\).  In the thesis, map 3 has target
\[
AR_1\times A^{-3}R_2\times AR_3\times A^{-1}R_4
\]
with the additional condition \(A\le (R_4/R_3)^{1/2}\), and its technical variant
has target \(AR_1\times A^{-3}R_2\times A\times B\) with
\(R_3<A<(R_3R_4)^{1/2}\) and \(AB=R_3R_4\).  These conditions exclude the symmetric
hard model with \(R_3\sim R_4\) and small stretch \(A=L^{1/5}\).  Thus no explicit
Guth-map counterexample to the theorem is currently known from the listed
Chapter 8 constructions.




Further consolidation of the averaged-filling route.

The non-absorbing averaged filling estimate
\[
\int_Q\Fill_R(Z_y)\,dy
\le C\left(R_1I+\frac{I^2}{S_1q}
+\frac{R_3R_4}{S_3S_4}F_0q\right)
\]
is false in this universal form.  Product counterexample:
\[
S=(1,1,L,L),\qquad R=(1,M,L,L),\qquad
f(u_1,u_2,u_3,u_4)=(u_1,u_2/M,u_3,u_4),
\]
with \(1<M<L\).  Then \(\Dil_2(f)=1\), the map has degree one, and for
\(F=(f_3,f_4)\) the central fibers are
\([0,1]\times[0,M]\times\{y_3,y_4\}\).  On these fibers
\(\lambda=\|d(f_1,f_2)|_{\ker dF}\|=1\), hence
\[
I\sim ML^2,\qquad q\sim L^2,\qquad F_0=S_1S_2S_3=L.
\]
A source calibration \(du_1\wedge du_2\wedge d\psi(u_3,u_4)\), with
\(\psi=0\) on the \((u_3,u_4)\)-boundary and \(\psi(y)\gtrsim L\), gives
\[
\Fill_R(Z_y)\gtrsim ML,\qquad \int_Q\Fill_R(Z_y)\,dy\gtrsim ML^3.
\]
The proposed right side is \(O(ML^2+M^2L^2+L^3)\), so \(M=L^{1/2}\) violates it
by \(L^{1/2}\).  This example has \(R_3R_4/S_3S_4=1\), so it does not attack the
final theorem in the hard regime, but it rules out any universal non-absorbing
AF estimate.

The natural replacement is the self-absorbing form
\[
A:=\int_Q\Fill_R(Z_y)\,dy
\le C\left(R_1I+\frac{I^2}{S_1q}\right)
 + C\frac{R_3R_4}{S_3S_4}A . \tag{SAF}
\]
In the hard regime \(R_3R_4\le \kappa S_3S_4\), the final term can be absorbed,
and the same algebra as before forces
\[
I\gtrsim qS_1(S_2S_3)^{1/2},
\]
hence by weighted coarea
\[
\Vol(R)\gtrsim S_1S_2^{1/2}S_3^{3/2}S_4.
\]
The product counterexample above is compatible with SAF because the coefficient
of \(A\) is order \(1\).

Thoughts on proving SAF.  A fiberwise filling inequality cannot work: central
product sheets have large filling caused by long normal travel while their
weighted \(G=(f_1,f_2)\)-energy is small.  The estimate must be applied to the
whole two-parameter complex of slices.  The heuristic is to discretize \(Q\) and
tighten the complex skeleton-by-skeleton as in Guth's proof of Estimate 1.  The
controlled part should contribute \(R_1I+I^2/(S_1q)\).  The cells for which the
tightening must use actual source normal area rather than target parameter area
should be charged not to the fixed lower target filling \(F_0q\) but to the
original average \(A\); their total proportion should be
\(R_3R_4/(S_3S_4)\).

One-dimensional slicing attempts do not currently close the gap.  For
\(H=(f_1,f_3,f_4)\), a central fiber maps to the \(x_2\)-segment in \(S\).
The target relative filling area of that segment is only \(\sim S_1S_2\), since
it can fill through the short \(x_1\)-direction.  Even if one could convert this
to a source length lower bound and integrate over parameters, the resulting
lower bound is at best of basic-volume scale and misses the factor
\((S_3/S_2)^{1/2}\).  Boundary ellipsoid inequality B3 gets the desired
\(\sqrt{S_2S_3}\) length scale because the circles are constrained to fill in
the boundary, where the short \(S_1\)-direction is unavailable.  An interior
argument needs an averaged/parametric way to reproduce this boundary obstruction
over an \(S_1\)-thick family, not just isolated level sets.

Guth's product-domain generalization near the end of arXiv:0710.0403 says that
for \(U\subset X^j\times Y^{l-j}\times Z^{n-l}\), a \(k\)-contracting
nonzero-degree map to \(S\) forces
\[
 |X|^{(l-j)/(k-j)}|Y|\gtrsim
 (S_1\cdots S_j)^{(l-j)/(k-j)}S_{j+1}\cdots S_l.
\]
With the standard product decomposition of \(R\) this recovers the monomial
estimates.  Attempts to choose a more favorable product by replacing the
\((x_3,x_4)\)-rectangle with a square of side \(\sqrt{R_3R_4}\) run into exactly
the need for a controlled projection/fiber-volume estimate; no such projection
is available for a highly eccentric rectangle without paying the lost aspect
factor.  Thus this published product estimate does not by itself imply the
desired mixed volume bound.



Further round-10 exploration: projection concentration and level-set routes.

For central slices \(Z_y=\langle [R,\partial R],F,y\rangle\), \(F=(f_3,f_4)\),
define
\[
 a(y)=\mathcal L^2(\pi_{12}(\operatorname{spt}Z_y)).
\]
A rigorous Fubini estimate is
\[
 \int_Q a(y)\,dy\le R_1R_2R_3R_4.
\]
Indeed \(u\in\pi_{12}(\operatorname{spt}Z_y)\) implies
\(y\in F(\{u\}\times[0,R_3]\times[0,R_4])\), whose area is at most \(R_3R_4\)
because \(J_2F\le1\) on the vertical rectangle.  In the hard regime this says
that the set of parameters for which the slice has large \((x_1,x_2)\)-base
projection has relative measure \(O(R_3R_4/(S_3S_4))\).  This exactly matches
the desired absorbing coefficient.

Tempting but currently unproved/false strengthening.  If one could prove a
good-slice inequality
\[
 \Fill_R(C)\le C R_1\Mass(C)
\]
for every relative 2-cycle \(C\) whose base projection omits a point (or has
area at most half of the base), then the theorem would follow immediately:
good central fibers would have mass at least \(S_1S_2S_3/R_1\), and coarea over
a \(1-O(\kappa)\) fraction of \(Q\) would give a volume bound stronger than the
desired one.  The obstruction is topological: retracting a punctured
\(R_1\times R_2\) base to its boundary as a pair forces a loop around the
puncture to wind once around the long outer boundary, and Guth's averaging
argument naturally produces an \(R_2\), not \(R_1\), loss.  The model
\(C=[0,R_1]\times[0,L]\) sitting at fixed values of the other coordinates shows
how boundary tracks can force a long annular sweep even when the two-dimensional
base projection has measure zero.  Thus a naive ``short horizontal retraction''
does not supply the missing estimate.

A weaker good-slice inequality with an \(R_2\) factor gives only
\[
 V_R\,R_2\gtrsim S_1S_2S_3^2S_4,
\]
the same scale as one of the double/boundary homotopy inequalities and still
compatible with the scaling test \(S=(1,1,L^6,L^6)\),
\(R=(L^{-1},L^4,L^5,L^6)\).  It therefore cannot prove the theorem.

Level-set attempt.  Slice by \(H=(f_1,f_3,f_4)\).  If for most
\((t,y_3,y_4)\) the preimage of the \(x_2\)-segment in the level
\(\{s_1=t,s_3=y_3,s_4=y_4\}\) had length \(\gtrsim (S_2S_3)^{1/2}\), then
\(\Dil_3(f)\le1\) and coarea would give exactly
\(V_R\gtrsim S_1S_2^{1/2}S_3^{3/2}S_4\).  The available filling obstruction for
that segment in the full target is only \(S_1S_2\), because it can fill through
the short \(s_1\)-direction.  The desired \(S_2S_3\) filling area is a boundary
or fixed-level obstruction, but the source level sets \(f_1=t\) have no uniform
isoperimetric inequality.  This route therefore also seems to require a
parametric argument that keeps track of level compatibility, not just isolated
curves.

Conclusion preserved for later work: the promising rigorous data are the
uniform target filling lower bound, the weighted mixed inequality
\(\lambda^2J_2F\le1\), and the projection average above.  The missing step is a
decoupling/tightening theorem that converts the unweighted small bad-projection
fraction into the self-absorbing term
\((R_3R_4/S_3S_4)\int_Q\Fill_R(Z_y)\,dy\), while charging the good part to
\(R_1I+I^2/(S_1|Q|)\) or to an equivalent quantity strong enough to force
\(I\gtrsim |Q|S_1(S_2S_3)^{1/2}\).



\section*{Round 11 algebra and strategy notes}

After the first-alternative reduction, the existing central-slice lower bound
is slightly stronger than previously emphasized.  Let
\(\alpha=R_1/S_1\) and
\[
T=S_1S_2^{1/2}S_3^{3/2}S_4,\qquad F_0=S_1S_2S_3,\qquad q\sim S_3S_4 .
\]
If \(R_3R_4\lesssim S_3S_4\), Proposition first gives
\(R_1R_2^2\gtrsim F_0\).  In the lower bound
\[
V_R\gtrsim q\min\{R_1R_2,F_0^{2/3},(R_1F_0)^{1/2}\},
\]
the first term is then at least \((R_1F_0)^{1/2}\), and the second is also at
least this because \(R_1\le S_1\le S_2\le S_3\).  Hence
\[
V_R\gtrsim \alpha^{1/2}T.
\]
Guth's \(j=1,k=2,l=4\) monomial estimate in the sharper variable \(\alpha\) gives
\[
V_R\gtrsim \alpha^{-2}S_1S_2S_3S_4.
\]
Therefore a genuine counterexample to the desired second alternative
\(V_R\gtrsim\kappa T\), under \(R_1\le \kappa S_1\), would have to satisfy
\[
\alpha\lesssim \kappa^2,\qquad {S_3\over S_2}\gtrsim \alpha^{-4}\kappa^{-2}.
\]
The unresolved range is thus not merely \(S_3/S_2\gg_\kappa1\), but the extreme
case where the first side is much smaller than the hypothesis requires and the
target aspect ratio is enormous.  This is precisely where the \(A^2/R_1\) branch
of the fiberwise filling profile loses the factor \(\alpha^{1/2}\).

Potential counterexample caveat checked this round: a tempting abstract
argument from Guth's area-expanding embeddings \(S\hookrightarrow R\) to a
degree-one pair map \(R\to S\) is not automatic.  The Pontryagin--Thom collapse
naturally gives a map \(R/\partial R\to S/\partial S\), but to lift it to an
actual map of pairs one must extend the inverse boundary map through the
complement as a map to \(\partial S\).  The boundary degree obstruction forces
the outer boundary map also to have degree one, so a constant collapse outside
the embedded copy is not a continuous pair map.  Any real counterexample must
therefore use explicit degree-one snake/pinching maps of pairs, not just an
opposite-direction embedding.

One-dimensional slicing check: slicing by \((f_1,f_3,f_4)\) gives preimage curves
mapping with degree one to \(x_2\)-segments.  In the full target those segments
fill through the short \(x_1\)-direction with area only \(\sim S_1S_2\), not
\(\sim S_2S_3\).  Thus the resulting curve-length lower bounds recover at most
basic-volume scale (or low-aspect partial ranges) and do not address the extreme
range above.  Any successful argument still appears to need a parametric
mechanism that prevents all \(x_2\)-segments, over an \(S_1\)-thick family, from
using the short \(x_1\)-escape independently.

Useful question for external help: Is there a published boundary or minimax
(Steenrod-square/cup-product) estimate for the two-parameter family of central
2-cycles \(Z_y=\langle[R],(f_3,f_4),y\rangle\) that gives a lower bound
\(\int_Q\int_{Z_y}\lambda^2\gtrsim |Q|S_1(S_2S_3)^{1/2}\) under
\(R_3R_4\ll S_3S_4\)?  Equivalently, can Guth's tightening construction be
localized to prove the self-absorbing estimate (SAF) recorded above?


\end{document}
